% ============================================================================
% 07_estudo_caso_pium_to.tex
% Estudo de caso operacional: foco real do Painel do Fogo (CENSIPAM) em
% Pium--TO, 12/05/2026. Validação independente do sistema de previsão,
% diagnóstico de descompasso entre dados de grade global e realidade
% local, e melhorias incorporadas ao pipeline em produção.
%
% Como usar:
%   Inserir como nova seção dentro do Capítulo 5 (Resultados), após a
%   Seção 5.8 "Validação operacional via aplicação web", como
%   subseção complementar de validação externa. Numeração sugerida:
%   5.9 "Validação cruzada com fonte externa independente: caso Pium--TO".
%   Dentro de "Implementação das melhorias": subsubseções M1--M4 (auditoria,
%   fusão INMET, retreino OOD, extensões API/auditoria alinhadas à literatura).
% ============================================================================

\section{Validação cruzada com fonte externa independente: caso Pium--TO}
\label{sec:res_estudo_caso_pium}

A validação operacional via \emph{smoke tests} controlados (Seção~\ref{sec:res_app}) demonstrou a coerência interna do sistema. Para uma validação externa mais exigente, foi escolhido um \emph{evento real} de incêndio confirmado por uma fonte oficial independente --- o \emph{Painel do Fogo} do Centro Gestor e Operacional do Sistema de Proteção da Amazônia (CENSIPAM)\footnote{\url{https://painel.censipam.gov.br/}, acessado em 12/05/2026 às 18:21 (UTC$-$3).} --- e submetido ao \emph{pipeline} em modo de produção. Esta seção apresenta o evento, as respostas do modelo, o diagnóstico das causas de divergência observada e as três correções incorporadas ao código em consequência da investigação.

% ----------------------------------------------------------------------------
\subsection{O evento}
\label{subsec:pium_evento}

O \emph{Painel do Fogo} listou um foco ativo identificado pelo \texttt{ID 6668060} em \texttt{(-10{,}5351; -49{,}8230)}, dentro do município de \textbf{Pium}, estado do \textbf{Tocantins}, em domínio fundiário misto (área privada e terra pública não categorizada) sobre cobertura de \emph{vegetação natural primária} (classificação TerraClass). Pium fica na ecorregião do Cerrado/Bico do Papagaio, dentro da Amazônia Legal. A Tabela~\ref{tab:pium_evento} reproduz as detecções oficiais.

\begin{table}[h]
\centering
\caption{Detecções no entorno de \texttt{(-10{,}5351; -49{,}8230)} --- evento \texttt{6668060} (\emph{Painel do Fogo}/CENSIPAM).}
\label{tab:pium_evento}
\small
\begin{tabular}{lcccc}
\toprule
\textbf{Data e hora (BR)} & \textbf{Área (ha)} & \textbf{Focos} & \textbf{Duração (h)} & \textbf{Propagação (ha/h)} \\
\midrule
12/05/2026 13:28 & 464,6 & 1  & \textbf{48,3} & 0,4 \\
10/05/2026 14:05 & 456,3 & 12 & 0,9  & 0,0 \\
10/05/2026 13:13 & 83,1  & 3  & 0,0  & 0,0 \\
\bottomrule
\end{tabular}
\end{table}

O status reportado pelo CENSIPAM era \emph{ativo pela última vez em 12/05/2026 14:49 (hora oficial do Brasil)}, com duração acumulada de 48,3~h --- compatível com vegetação ressecada queimando em regime sustentado.

% ----------------------------------------------------------------------------
\subsection{Resposta do sistema: cinco janelas temporais}
\label{subsec:pium_resposta}

A \emph{API} \texttt{/api/predict} foi consultada em cinco janelas centradas no evento, usando o modelo final \texttt{ensemble\_stacking\_gbm}. Em todas elas o sistema previu \textbf{\emph{Baixo}} com probabilidade entre 81 e 88\%, conforme Tabela~\ref{tab:pium_respostas_pos_fix} (\emph{já após} a correção dos \emph{bugs} identificados na Subseção~\ref{subsec:pium_correcoes}).

\begin{table}[h]
\centering
\caption{Respostas do \texttt{ensemble\_stacking\_gbm} para Pium--TO em diferentes janelas (modelo \emph{argmax}, após correções).}
\label{tab:pium_respostas_pos_fix}
\small
\begin{tabular}{lccccc}
\toprule
\textbf{Janela consultada}     & \textbf{Risco} & \textbf{P(Baixo)} & \textbf{P(Mod.)} & \textbf{P(M.~Alto)} & \textbf{FRP máx.~5d} \\
\midrule
12/05/2026 14:00 & Baixo & 85,9\% & 9,5\%  & 4,6\% & \SI{16,6}{\mega\watt} (16~detec.) \\
11/05/2026 12:00 & Baixo & 85,8\% & 9,7\%  & 4,5\% & \SI{22,0}{\mega\watt} (19~detec.) \\
10/05/2026 13:00 & Baixo & 87,1\% & 8,6\%  & 4,3\% & \SI{22,0}{\mega\watt} (21~detec.) \\
09/05/2026 12:00 & Baixo & 86,6\% & 8,8\%  & 4,6\% & \SI{22,0}{\mega\watt} (26~detec.) \\
05/05/2026 12:00 & Baixo & 81,4\% & 14,3\% & 4,3\% & \SI{58,4}{\mega\watt} (32~detec.) \\
\bottomrule
\end{tabular}
\end{table}

A divergência entre a realidade (foco confirmado e duração de 48,3~h) e a saída do modelo (\emph{Baixo} consistente em todas as janelas) constituiu o ponto de partida para a investigação descrita a seguir.

% ----------------------------------------------------------------------------
\subsection{Diagnóstico das causas}
\label{subsec:pium_diagnostico}

A inspeção do vetor de \emph{features} entregue ao classificador identificou \textbf{três causas convergentes}, descritas em ordem decrescente de impacto.

% ............................................................................
\paragraph{Causa 1 (corrigida): integração com NASA FIRMS retornando \texttt{HTTP 400}.}

O cliente da NASA FIRMS (\texttt{scripts/frp\_api.py}) estava emitindo requisições com \texttt{days\_back=7} para o \emph{endpoint} de \emph{area} no \emph{dataset} VIIRS\_SNPP\_NRT. A resposta do servidor, em todas as datas testadas, foi:

\begin{quote}
\texttt{HTTP/1.1 400 Bad Request\\Invalid day range. Expects [1..5].}
\end{quote}

Verificação direta do \emph{status} da \texttt{MAP\_KEY} confirmou que a chave estava válida (\texttt{transaction\_limit=5000}, \texttt{current\_transactions=0}), demonstrando que o erro era \emph{parametrização}, não \emph{credenciais}. A API NASA FIRMS NRT passou a aceitar no máximo 5~dias por requisição\footnote{Documentação oficial: \url{https://firms.modaps.eosdis.nasa.gov/api/area/}, consultada em 12/05/2026.}.

Em consequência, as \emph{features} \texttt{Media\_FRP\_Ultimos\_7\_Dias}, \texttt{Max\_FRP\_Ultimos\_7\_Dias} e \texttt{Dias\_Desde\_Ultimo\_Incendio} estavam recebendo seus valores de \emph{fallback} (zeros e 246--365~dias), e o modelo ``não enxergava'' o fogo recente em curso. Após a correção (\(\texttt{days\_back} \in [1,5]\) e propagação de \texttt{reference\_date}), o sistema passou a recuperar até 32 detecções VIIRS no raio de 10~km, com FRP máximo de \SI{58,4}{\mega\watt}.

% ............................................................................
\paragraph{Causa 2 (corrigida): \emph{lookup} Tier~1 degradando precocemente para \texttt{(Estado, Mês)}.}

A primeira versão do \emph{lookup} espaço-sazonal (\texttt{scripts/feature\_lookup.py}, Seção~\ref{subsec:metod_features_tier1}) buscava medianas em uma janela ${3{\times}3}$ células ($\approx 25$~km no equador, dada a resolução \SI{0,25}{\degree}). Para o ponto consultado em Pium--TO, mês~5, essa janela retornou vazia e o sistema degradou imediatamente para a granularidade \texttt{estado\_mes} (mediana de todo o Tocantins em maio). O efeito foi:

\begin{itemize}
  \item \texttt{Precipitacao\_acum\_30/90/180/365d} = \SI{0,0}{mm} (mediana estadual);
  \item \texttt{SPI\_1m} = \texttt{SPI\_3m} = \texttt{SPI\_6m} = 0,00 (climatologia neutra);
  \item \texttt{Dias\_Secos\_90d} = 2 (mediana estadual, não local);
  \item \texttt{Media\_FRP\_Celula\_30d} = \SI{10,65}{\mega\watt} (não específico).
\end{itemize}

Isto é, \emph{toda a informação espacial fina foi apagada}. A cascata de busca foi ampliada para
\(3{\times}3 \to 5{\times}5 \to \text{mês adjacente} \to (\text{Estado, Mês}) \to (\text{Mês global})\),
elevando a granularidade efetiva para \texttt{celula\_ext} (anel externo, $\approx 55$~km) no caso Pium--TO. Os valores de \texttt{Precipitacao\_acum\_*d} subiram para \SI{0,5}{mm}, \texttt{Dias\_Secos\_90d} para 1 e \texttt{Media\_FRP\_Celula\_30d} para \SI{8,2}{\mega\watt} --- valores ainda baixos, mas vindos de células geograficamente próximas (Tabela~\ref{tab:pium_tier1_pre_pos}).

\begin{table}[h]
\centering
\caption{Subconjunto das \emph{features} Tier 1 antes e depois da expansão da cascata espacial.}
\label{tab:pium_tier1_pre_pos}
\small
\begin{tabular}{lcc}
\toprule
\textbf{Feature} & \textbf{Cascata \texttt{3x3} (\texttt{estado\_mes})} & \textbf{Cascata estendida (\texttt{celula\_ext})} \\
\midrule
\texttt{Precipitacao\_acum\_30d}      & 0,0 mm  & 0,5 mm \\
\texttt{Precipitacao\_acum\_180d}     & 0,0 mm  & 0,5 mm \\
\texttt{Precipitacao\_acum\_365d}     & 1,3 mm  & 0,5 mm \\
\texttt{SPI\_3m}                       & 0,00    & 0,00 \\
\texttt{Dias\_Secos\_90d}              & 2       & 1 \\
\texttt{Incendios\_Ultimos\_90\_Dias}  & 2       & 1 \\
\texttt{Media\_FRP\_Celula\_30d}       & 10,65 MW & 8,2 MW \\
\bottomrule
\end{tabular}
\end{table}

% ............................................................................
\paragraph{Causa 3 (não corrigível em produção, documentada): descompasso entre clima de grade global e clima local em zona de transição.}

Após as duas correções acima, a previsão para 12/05/2026 \emph{aumentou} a confiança em \emph{Baixo}, de \SI{82,7}{\percent} para \SI{85,9}{\percent}. A investigação detalhada do vetor de \emph{features} revelou que a NASA POWER MERRA-2 (grade $\approx 50$~km) reportou para a janela de 28~dias anterior a 12/05/2026:

\begin{quote}
\texttt{Precipitacao\_ma7 = 0,12 mm/dia}, \texttt{DiaSemChuva\_ma7 = 1,3 dia}, \texttt{RH2M = 64,6\%}.
\end{quote}

Esses valores caracterizam regime \emph{ainda chuvoso} (chuvas frequentes mas leves). Para o modelo, classificar como \emph{Baixo} foi a decisão coerente com os \emph{inputs} climáticos recebidos. O problema é que esses \emph{inputs} não corresponderam à realidade local: o foco real ardia há mais de 48~h.

Para confirmar que o modelo \emph{não} é cego sazonalmente, foi executado um experimento \emph{contrafactual} no mesmo ponto \texttt{(-10{,}5351; -49{,}8230)}, variando a data (Tabela~\ref{tab:pium_contrafactual}).

\begin{table}[h]
\centering
\caption{Experimento contrafactual --- mesmo ponto, climatologia distinta capturada pela NASA POWER.}
\label{tab:pium_contrafactual}
\small
\begin{tabular}{lcccccc}
\toprule
\textbf{Data} & \textbf{KBDI \emph{proxy}} & \textbf{Prec.~ma7} & \textbf{DSC~ma7} & \textbf{FRP medido} & \textbf{P(M.~Alto)} & \textbf{Risco} \\
\midrule
15/08/2024 & 424,1 & 0,02 & 16,0 dias  & \SI{0}{\mega\watt} & 99,0\% & \textbf{Muito Alto} \\
15/09/2024 & 712,5 & 0,00 & 25,0 dias  & \SI{0}{\mega\watt} & 99,6\% & \textbf{Muito Alto} \\
15/05/2024 & 148,0 & 0,51 & 8,6 dias   & \SI{0}{\mega\watt} & 47,3\% & \textbf{Muito Alto} \\
12/05/2026 & 29,8  & 0,12 & 1,3 dia    & \SI{16,6}{\mega\watt} & 4,6\% & Baixo \\
\bottomrule
\end{tabular}
\end{table}

Três conclusões emergem:

\begin{enumerate}
  \item O modelo \emph{responde corretamente} ao sinal climático: 15/05/2024 (mesmo dia e mesmo ponto, mas em ano de estiagem mais marcada) recebeu classificação \emph{Muito Alto} com confiança \SI{47,3}{\percent}, mesmo sem FRP detectado.
  \item A previsão \emph{Muito Alto} se intensifica monotonicamente conforme cresce o estresse hídrico medido pela NASA POWER (KBDI~$424 \to 712$, DSC$_\text{ma7}$~$16 \to 25$~dias).
  \item Em 12/05/2026 (caso real), os indicadores climáticos da grade MERRA-2 estavam \emph{suaves} (KBDI=29,8; DSC=1,3~dia; Prec\_ma7=0,12~mm/dia) --- e o modelo refletiu isso. \textbf{O erro está no \emph{input}, não na decisão}.
\end{enumerate}

A interpretação operacional é direta: a célula MERRA-2 que cobre Pium--TO em maio de 2026 \emph{captou chuvas localizadas/dispersas} suficientes para elevar a frequência de dias com precipitação, mas não a quantidade total --- regime ``chuvas leves e frequentes''. A vegetação alvo, em escala municipal, continuou em condição de queima. Esse é um caso clássico de \emph{representativity gap} entre dados de grade e fenômenos locais~\cite{seager2015climatology}, particularmente relevante em zonas de transição Cerrado--Amazônia onde a heterogeneidade climática intra-célula é grande.

% ............................................................................
\paragraph{Hipótese descartada: sobrescrever contagem de incêndios com observação FIRMS.}

Uma tentativa intermediária consistiu em substituir as \emph{features} \texttt{Incendios\_Ultimos\_7\_Dias} e \texttt{Incendios\_Ultimos\_30\_Dias} pela contagem real de detecções FIRMS (16--32~focos no caso). \emph{Surpreendentemente}, isso \textbf{aumentou} a confiança em \emph{Baixo} (de \SI{82,7}{\percent} para \SI{86,0}{\percent}). A causa é \emph{distribution shift}: no \emph{dataset} de treino, essas \emph{features} para Tocantins--Maio têm mediana zero; ao injetar valores de duas ordens de magnitude maiores, o exemplo é colocado \emph{fora da distribuição} (\emph{out-of-distribution}, OOD), e o classificador responde de modo errático. A sobrescrita foi revertida; mantêm-se atualizadas apenas \texttt{FRP\_max\_7d}, \texttt{Media\_FRP\_7d} e \texttt{Dias\_Desde\_Ultimo\_Incendio}, que no \emph{dataset} de treino eram observações individuais (não medianas).

% ----------------------------------------------------------------------------
\subsection{Correções incorporadas e seu impacto}
\label{subsec:pium_correcoes}

A Tabela~\ref{tab:pium_correcoes} sintetiza as três correções introduzidas no \emph{pipeline} em consequência do estudo de caso e o impacto medido para a janela 12/05/2026 14:00.

\begin{table}[h]
\centering
\caption{Correções incorporadas e impacto medido (Pium--TO, 12/05/2026 14:00).}
\label{tab:pium_correcoes}
\small
\begin{tabular}{p{4.6cm}p{5.2cm}p{2.4cm}p{1.6cm}}
\toprule
\textbf{Correção} & \textbf{Mudança técnica} & \textbf{Granularidade Tier 1} & \textbf{P(Baixo)} \\
\midrule
\emph{Baseline} (antes do caso) & --- & \texttt{estado\_mes} & 82,0\% \\
(C1) NASA FIRMS \texttt{HTTP 400} & \texttt{days\_back $\in [1,5]$}; aceita \texttt{reference\_date}; URL com sufixo \texttt{/YYYY-MM-DD} & \texttt{estado\_mes} & 82,7\% \\
(C2) Cascata Tier 1 ampliada & $3{\times}3 \to 5{\times}5 \to$ mês adj. & \texttt{celula\_ext} & 85,9\% \\
(C3) Inc.~Últ.~$N$~dias mantida como mediana de treino & Apenas \texttt{FRP\_*} e \texttt{Dias\_Desde\_Ultimo\_Incendio} são atualizados em tempo real & \texttt{celula\_ext} & 85,9\% \\
\bottomrule
\end{tabular}
\end{table}

As correções \textbf{(C1)} e \textbf{(C2)} são melhorias técnicas inequívocas. A correção \textbf{(C3)} é uma \emph{decisão de engenharia conservadora}: preferir o sinal aprendido pelo modelo (mediana sazonal Estado--Mês, que é o que ele viu durante o treino) a inputs reais que o colocam OOD. A elevação aparentemente paradoxal de P(Baixo) após as correções é coerente com a Causa~3: o modelo passou a \emph{ver melhor} um cenário que --- para a grade MERRA-2 --- parecia chuvoso.

% ----------------------------------------------------------------------------
\subsection{Discussão: limites do paradigma de dados de grade e caminhos para mitigação}
\label{subsec:pium_discussao}

O caso Pium--TO 2026 expõe uma tensão estrutural do nosso \emph{pipeline} (e, mais amplamente, de qualquer sistema de previsão de fogo apoiado em reanálise climática global): a granularidade espacial dos \emph{inputs} climáticos limita a resolução máxima do \emph{output}. NASA POWER MERRA-2~\cite{nasapower} entrega cobertura global e estabilidade temporal, mas com células de $\approx \SI{50}{km}$ que suavizam microclimas. Em regiões homogêneas (centro do Amazonas), a perda de resolução é benigna; em regiões de transição (Cerrado--Amazônia, Pantanal, Caatinga), a perda pode mascarar condições propícias a fogo.

Três caminhos são discutidos como mitigação no Capítulo~\ref{cap:conclusao}:

\begin{enumerate}
  \item \textbf{Fusão MERRA-2 + INMET}~\cite{inmetnormais}: utilizar precipitação de estação meteorológica quando houver estação INMET dentro de $\sim 50$~km do ponto consultado (ou, em modo hierárquico configurável, até $\sim 180$~km com rótulo explícito de \texttt{inmet\_representatividade}, ver Subsubseção~M4). Para dados em tempo real, a API WIS2/OGC do INMET (\texttt{wis2bra.inmet.gov.br/oapi}) expõe observações horárias de 358+ estações sinópticas sem necessidade de token, com histórico de $\sim$90 dias; para histórico mais profundo, os ZIPs anuais oficiais cobrem 565+ estações automáticas desde 2000. Trade-off: cobertura esparsa em zonas remotas, exige rotina de \emph{nearest station}, \emph{quality control} e transparência sobre a distância da estação.
  \item \textbf{Aumento de exemplos OOD no treino}: oversampling estratégico de focos confirmados FIRMS/INPE em zonas de transição (Cerrado em maio, Pantanal em junho), reduzindo o \emph{prior} de raridade que o modelo aprendeu para esses pares \texttt{(Estado, Mês)}. Trade-off: risco de degradar a performance nas regiões de regime mais homogêneo.
  \item \textbf{Auditoria operacional contínua}: \emph{logging} estruturado de \texttt{(coordenada, data, risco predito, FRP detectado depois)} para calcular $\text{F}_1$ operacional \emph{rolling} em janelas móveis e detectar deriva de modelo em produção --- prática alinhada com o paradigma de \emph{ML observability}.
\end{enumerate}

% ----------------------------------------------------------------------------
\subsection{Implementação das três melhorias propostas e validação empírica}
\label{subsec:pium_melhorias}

As três melhorias listadas na Seção~\ref{subsec:pium_discussao} foram implementadas e validadas em código de produção. Esta subseção registra o resultado empírico de cada uma.

\subsubsection*{M1. Auditoria operacional contínua}

Foi adicionado o módulo \texttt{scripts/audit\_log.py} que escreve, a cada chamada de \texttt{/api/predict}, uma linha JSONL em \texttt{logs/auditoria\_predicoes.jsonl} contendo: \texttt{prediction\_id}, coordenadas, \texttt{ref\_data}, classe predita, probabilidades, modelo, 15 \emph{features} chave, fonte climática (com detalhe INMET se aplicável) e granularidade Tier 1. O log é \emph{append-only} com rotação automática a 50~MB e tolerante a falhas (a predição segue caso a gravação falhe).

O script \texttt{scripts/auditar\_predicoes.py} consome esse JSONL, consulta a NASA FIRMS NRT na janela $[t, t+H]$ para cada predição com idade $\geq d_{\min}$ (defaults $H=3$~d, $d_{\min}=3$~d, raio 5~km), define $y_{\text{true}} \in \{0,1\}$ (foco confirmado) e $y_{\text{pred}} \in \{0,1\}$ (Moderado/Muito Alto $\to$ 1) e calcula precision, recall, $\text{F}_1$ e acurácia globais e em janelas \emph{rolling} de 7, 14 e 30 dias. Predições mais antigas que o histórico do NRT ($\sim 60$~dias) aparecem no CSV mas são automaticamente excluídas das métricas.

\subsubsection*{M2. Fusão NASA POWER + INMET (duas fontes complementares)}

A integração foi implementada com \emph{duas} fontes INMET, selecionadas automaticamente conforme a idade da consulta. Esta arquitetura híbrida foi adotada porque a sondagem ativa em 12/05/2026 mostrou que (i) a API horária tradicional \texttt{apitempo.inmet.gov.br/estacao/...} retorna \texttt{status 204} com frequência e não é confiável; (ii) o ZIP histórico é completo mas o do ano corrente é parcial (publicado mensalmente, em 12/05/2026 cobre só janeiro--fevereiro/2026); (iii) o INMET passou a publicar dados horários sem token via a \textbf{WIS2 (OGC API Features)}~\cite{inmetnormais}, ainda pouco documentada em literatura brasileira.

\textbf{Fonte (1)~--~WIS2/OGC API Features} (\texttt{http://wis2bra.inmet.gov.br/oapi}): expõe a variável \texttt{total\_precipitation\_or\_total\_water\_equivalent} (em \texttt{kg\,m$^{-2}$}~$\equiv$~mm/h) na coleção \texttt{urn:wmo:md:br-inmet:synop} para 358 estações sinópticas WMO brasileiras que publicaram chuva nos últimos 7 dias, com histórico de $\sim$90 dias e \emph{sem necessidade de token}. A cobertura geográfica é menor que a do ZIP (só estações sinópticas oficiais, não todas as automáticas), mas a latência é de minutos. O endpoint possui um detalhe: o filtro \texttt{wigos\_station\_identifier} na \emph{querystring} é \emph{ignorado} (verificado empiricamente em 12/05/2026), e a filtragem precisa ser feita \emph{client-side} após restringir o \texttt{bbox} a $\sim$5~km em torno da estação para evitar baixar 58~mil observações misturadas.

\textbf{Fonte (2)~--~ZIPs anuais} (\texttt{portal.inmet.gov.br/uploads/dadoshistoricos/\{ANO\}.zip}, $\sim$100~MB/ano): cobertura completa de 565+ estações automáticas desde 2000 (168 na Amazônia Legal). Cache local em \texttt{.cache\_inmet/\{ANO\}/\{COD\}.csv} evita re-download. O encoding é \texttt{latin-1}, separador \texttt{;}, decimal \texttt{,}.

A função \texttt{get\_inmet\_fusion\_data} roteia automaticamente: se \texttt{reference\_date} está nos últimos 90 dias \emph{e} existe estação WIS2 com chuva a $\leq 50$~km, usa WIS2 (sem download, $\sim$1~s); senão, tenta o ZIP anual. Se nenhuma das duas funciona, retorna \texttt{None} e o \texttt{climate\_api} cai graciosamente no MERRA-2.

A fusão em \texttt{climate\_api.get\_climate\_data} prioriza NASA POWER (sempre disponível) e, quando o INMET retorna dados (WIS2 ou ZIP), \emph{sobrescreve} \texttt{precipitacao}, \texttt{prec\_ma7}, \texttt{dias\_sem\_chuva} e \texttt{diasem\_ma7} pelos valores INMET. Os valores MERRA-2 originais são preservados em chaves \texttt{*\_merra} para auditoria, e a resposta da API expõe \texttt{fonte\_dados} $\in$ \{\texttt{nasa\_power}, \texttt{inmet\_fundido\_nasa\_power}\}, além de \texttt{estacao\_inmet}, \texttt{distancia\_estacao\_inmet\_km}, \texttt{cobertura\_inmet\_pct}. A Tabela~\ref{tab:pium_inmet} apresenta a validação empírica nos quatro cenários testados.

\begin{table}[h]
\centering
\caption{Validação da fusão INMET com duas fontes (WIS2 real-time e ZIP histórico).}
\label{tab:pium_inmet}
\small
\begin{tabular}{p{3cm}cp{2.6cm}cccc}
\toprule
\textbf{Caso} & \textbf{Fonte} & \textbf{Estação} & \textbf{Cob.} & \textbf{MERRA-2 DSC} & \textbf{INMET DSC} & \textbf{Risco} \\
\midrule
Pium-TO 15/08/2024 & ZIP & A055 Lagoa da Confusão (32,7~km) & 89\% & 16 d & \textbf{7 d} & Muito Alto (97,3\%) \\
Brasília 10/08/2024 & ZIP & A001 Brasília (1,1~km) & 89\% & 25 d & \textbf{7 d} & Muito Alto (67,9\%) \\
Pium-TO 12/05/2026 & WIS2 + hier.\textsuperscript{$\dagger$} & FORMOSO ($\sim$152~km, \texttt{baixa}) & --- & 1,3 d & \textbf{7 d} & Baixo (85,9\%); metadados na API \\
\textbf{Brasília 12/05/2026 (D-0)} & \textbf{WIS2} & \textbf{A001 (1,1~km)} & \textbf{88\%} & \textbf{1,86 d} & \textbf{7 d} & \textbf{Baixo 66,0\%} (era $\sim$99\%) \\
\bottomrule
\end{tabular}\\
\footnotesize
\textsuperscript{$\dagger$} O ZIP do ano corrente é parcial; em maio/2026 o arquivo cobre só janeiro--fevereiro/2026. A estação WIS2 mais próxima de Pium-TO está a $\sim$152~km (FORMOSO DO ARAGUAIA-TO) --- fora do raio inicial de 50~km, mas coberta pela \textbf{busca hierárquica} (até 180~km) com \texttt{inmet\_representatividade = baixa}. O classificador em produção ainda usa o vetor de \emph{features} do retreino anterior; a linha da tabela documenta o contraste MERRA-2 vs.\ INMET na camada de fusão e na API.
\end{table}

O caso \textbf{Brasília 12/05/2026 (D-0)} é a primeira demonstração do fluxo de fusão \emph{real-time end-to-end}: às 19h26 BRT do dia da consulta, a WIS2 detectou que o MERRA-2 reportava um microclima úmido espúrio (DSC$_{\text{ma7}}=1{,}86$~d) enquanto a estação INMET A001 a 1,1~km do ponto consultado registrava 7 dias secos consecutivos. A correção foi propagada para o modelo e derrubou a confiança da classe Baixo de aproximadamente 99\% para 66\% --- sinal claro de que, com mais sub-amostragem regional e fusão mais agressiva, o modelo passaria a alertar.

A fusão demonstrou \emph{capacidade de correção} robusta tanto retrospectivamente (Pium-TO e Brasília em agosto/2024, via ZIP) quanto em tempo real (Brasília 12/05/2026 via WIS2). Para zonas remotas do Amazonas ou Sul do Pará, nenhuma estação está dentro do raio de 50~km e o sistema continua usando MERRA-2, como antes.

\subsubsection*{M3. Retreino com oversampling OOD em Cerrado-transição}

A análise da distribuição de focos no \texttt{base\_de\_dados\_enriquecido.csv} (924\,306 linhas, 547\,845 com FRP $> 0 = 58{,}7\%$) identificou \textbf{oito pares (Estado, Mês) sub-representados} em estados de transição Cerrado-Amazônia (TO, MA, MT, PA) nos meses de transição chuvoso--seco (maio--julho), incluindo o par \texttt{(TOCANTINS, 5)} relevante para o caso Pium-TO. A Tabela~\ref{tab:pium_ood_alvos} resume a calibração de oversampling.

\begin{table}[h]
\centering
\caption{Pares (Estado, Mês) selecionados para oversampling OOD.}
\label{tab:pium_ood_alvos}
\small
\begin{tabular}{lccccc}
\toprule
\textbf{Estado} & \textbf{Mês} & \textbf{$n_\text{total}$} & \textbf{$n_\text{focos}$} & \textbf{Razão vs.~média estado} & \textbf{Fator de oversampling} \\
\midrule
PARÁ & 5 & 1.293 & 885 & 0,05 & 5,0$\times$ \\
MARANHÃO & 5 & 158 & 115 & 0,06 & 5,0$\times$ \\
PARÁ & 6 & 4.442 & 2.922 & 0,18 & 5,0$\times$ \\
MARANHÃO & 6 & 617 & 383 & 0,20 & 5,0$\times$ \\
\textbf{TOCANTINS} & \textbf{5} & \textbf{75} & \textbf{50} & \textbf{0,29} & \textbf{3,5$\times$} \\
MARANHÃO & 7 & 1.536 & 817 & 0,42 & 2,4$\times$ \\
PARÁ & 7 & 17.262 & 10.139 & 0,62 & 1,6$\times$ \\
TOCANTINS & 6 & 197 & 118 & 0,68 & 1,5$\times$ \\
\bottomrule
\end{tabular}
\end{table}

O oversampling é aplicado \emph{apenas no conjunto de treino} (\emph{split} estratificado, \texttt{random\_state=42}, mesmo do \emph{baseline}; teste com 184\,862 amostras intactas). Cada linha alvo é replicada $\lfloor \text{fator}-1 \rfloor$ vezes com \emph{jitter} gaussiano em Latitude/Longitude ($\sigma \approx 2{,}2$~km), Hora ($\sigma=1$~h) e FRP ($\sigma = 5\%$ do valor) --- para evitar duplicação exata sem alterar a semântica do par. O dataset aumentado teve +35\,472 linhas no treino (+4,8\%).

A Tabela~\ref{tab:pium_ood_metricas} apresenta os resultados para a versão calibratória (\texttt{logistic\_regression\_balanced}, 6,2~min de treino, mesma calibração isotônica do \emph{baseline}). O ganho é consistente em todas as classes, com destaque para Moderado (a classe mais difícil).

\begin{table}[h]
\centering
\caption{Comparativo OOD vs.~baseline no hold-out (n=184\,862) e revalidação do caso Pium-TO.}
\label{tab:pium_ood_metricas}
\small
\begin{tabular}{lccc}
\toprule
\textbf{Métrica} & \textbf{Baseline} & \textbf{Com OOD oversampling} & \textbf{$\Delta$ (p.p.)} \\
\midrule
Acurácia & 64,58\% & 65,75\% & \textbf{+1,17} \\
$\text{F}_1$-macro & 0,6044 & 0,6168 & \textbf{+1,25} \\
$\text{F}_1$-Baixo & 0,7142 & 0,7243 & +1,00 \\
$\text{F}_1$-Moderado & 0,3643 & 0,3793 & \textbf{+1,50} \\
$\text{F}_1$-Muito Alto & 0,7346 & 0,7470 & +1,24 \\
\midrule
\multicolumn{4}{l}{\emph{Revalidação Pium-TO --- P(Muito Alto)}} \\
12/05/2026 (foco real, MERRA-2 chuvoso) & 3,3\% & \textbf{6,6\%} & \textbf{\bm{$\times 2$}} \\
15/05/2024 (mesmo dia, MERRA-2 seco) & 23,1\% & \textbf{32,4\%} & +9,3 \\
15/08/2024 (pico de seca) & 66,2\% & \textbf{72,5\%} & +6,3 \\
\bottomrule
\end{tabular}
\end{table}

A classe final continua sendo ``Baixo'' para 12/05/2026 (predição $\neq$ realidade), mas a probabilidade de ``Muito Alto'' \emph{dobrou} (3,3\% $\to$ 6,6\%), e para o contrafactual com clima seco (15/05/2024) subiu 9,3~p.p. O modelo OOD ficou \emph{mensuravelmente} mais sensível a cenários de transição, exatamente como pretendido. O retreino do ensemble final \texttt{ensemble\_stacking\_gbm\_ood} ($\sim 90$~min de CPU) está parametrizado e disponível via:

\begin{center}
\texttt{python scripts/retreinar\_com\_oversample.py --modelo ensemble\_stacking\_gbm}
\end{center}

\noindent ele segue o mesmo \emph{split} estratificado e adicionalmente reporta um \emph{hold-out} restrito a Cerrado-transição (estados TO/MA/PA/MT/BA/GO/DF/PI em meses 5--7), que captura especificamente o tipo de cenário do caso Pium-TO.

\subsubsection*{M4. Extensões pós-literatura: hierarquia INMET, \emph{fire weather} na API e incerteza operacional}

Para alinhar o painel à literatura recente sobre \emph{fire weather}, fusão estação--grade e suporte à decisão sob ambiguidade --- \textbf{sem} alterar, por ora, as colunas do \texttt{ColumnTransformer} serializado (evita \emph{train/serve skew} até retreino documentado) --- foram acrescentados:

\begin{enumerate}
  \item \textbf{Busca INMET hierárquica} (\texttt{INMET\_EXTENDED\_MAX\_DISTANCE\_KM} em \texttt{scripts/config.py}; \texttt{get\_inmet\_fusion\_data(..., hierarchical=True)} em \texttt{scripts/inmet\_api.py}): se não há estação com cobertura dentro de 50~km, o sistema tenta até 180~km e preenche \texttt{inmet\_representatividade} $\in$ \{\texttt{alta}, \texttt{moderada}, \texttt{baixa}\} e \texttt{inmet\_busca\_raio\_km}. Para Pium--TO em 12/05/2026 isso recupera FORMOSO DO ARAGUAIA ($\sim$152~km, representatividade \texttt{baixa}), confrontando explicitamente o MERRA-2 úmido com medição in situ distante --- matéria direta de discussão sobre extrapolação espacial no Capítulo~\ref{cap:conclusao}.
  \item \textbf{Vento na mesma estação da precipitação} quando a fonte é WIS2 (\texttt{vento\_inmet\_ms\_ma7} em \texttt{scripts/climate\_api.py}), em linha com estudos que condicionam \emph{fire weather} a vento observado junto da chuva.
  \item \textbf{NASA POWER estendido} na requisição diária (\texttt{WS2M}, \texttt{T2M\_MAX} em \texttt{scripts/nasa\_power\_realtime.py}): médias móveis de 7~dias expostas ao painel (\texttt{ws2m\_ma7\_ms}, \texttt{t2m\_max\_ma7\_c}) para interpretabilidade e futuras ablações.
  \item \textbf{Proxy \texttt{FWI\_fire\_weather\_proxy}} e objeto \texttt{incerteza\_operacional} (\texttt{scripts/operational\_uncertainty.py}): indicador escalar combinando seca na semana, déficit de umidade relativa e vento; \emph{gap} entre as duas classes mais prováveis --- narrativa de suporte à decisão sob classes próximas, \emph{sem} alegar intervalos preditivos conformais calibrados em conjunto dedicado.
  \item \textbf{FNR/FPR na auditoria FIRMS} (\texttt{scripts/auditar\_predicoes.py}): o binário operacional alerta/não alerta passa a reportar taxa de falsos negativos e falsos positivos, frequentemente omitida quando só se divulga $\text{F}_1$.
\end{enumerate}

% ----------------------------------------------------------------------------
\subsection{Síntese}
\label{subsec:pium_sintese}

O estudo de caso Pium--TO 12/05/2026 contribuiu para o trabalho em três dimensões:

\begin{itemize}
  \item \textbf{Operacionalmente}: identificou e corrigiu dois \emph{bugs} silenciosos no \emph{pipeline} em produção (NASA FIRMS 400 e cascata Tier 1 prematura), e definiu uma política de \emph{feature fusion} entre observação NRT e medianas de treino (Tabela~\ref{tab:pium_correcoes}).
  \item \textbf{Cientificamente}: documentou empiricamente, com fonte externa independente, o limite imposto pela resolução da reanálise climática global sobre a previsão de fogo em zonas de transição, e estabeleceu o referencial \emph{contrafactual} (Tabela~\ref{tab:pium_contrafactual}) que mostra que o modelo \emph{não} sofre de viés sazonal cego --- ele responde corretamente aos \emph{inputs} que recebe.
  \item \textbf{Em melhorias entregues}: as três frentes apontadas como mitigação (Seção~\ref{subsec:pium_discussao}) foram implementadas e validadas (Seção~\ref{subsec:pium_melhorias}): auditoria operacional contínua (M1); fusão INMET com WIS2/ZIP e raio hierárquico (M2, Tabela~\ref{tab:pium_inmet}); retreino com oversampling OOD (M3, Tabela~\ref{tab:pium_ood_metricas}). A subsubseção~M4 acrescenta camada de interpretabilidade e governança (\emph{fire weather} na API, incerteza operacional heurística, FNR/FPR) alinhada ao estado da arte recente, deliberadamente \emph{fora} do vetor de treino até retreino controlado.
\end{itemize}

O caso completa a triangulação da validação do sistema: (i) \emph{hold-out} estratificado (Tabela~\ref{tab:res_comparativo}), (ii) validação temporal \emph{rolling-origin} (Tabela~\ref{tab:res_temporal_stacking}), (iii) \emph{smoke tests} operacionais controlados (Tabela~\ref{tab:res_smoke}) e (iv) confronto com fonte externa independente (esta seção). As quatro evidências, em conjunto, suportam a defensabilidade do modelo final \texttt{ensemble\_stacking\_gbm} para o uso aplicado proposto, com as limitações documentadas honestamente no Capítulo~\ref{cap:conclusao}.

% NOTA: ajustar a chave \cite{inmetnormais} se ainda não estiver em references.bib.
